% sections/ndss_04_measurement.tex  — NDSS 2027

\subsection{LSCPS: Static Capability Extraction}

\lscps (Layered Structured Capability Pattern Search) is a
rule-based static analysis pipeline that scans all skill artifact
files---Python/JS/shell code, SKILL.md instructions,
dependency manifests, and agent config files---for 300+
regex patterns organized into 14 security-sensitive capability families.
Patterns operate on parsed AST nodes where available
(Python \texttt{ast}, \texttt{esprima}) and raw text for shell scripts.
\lscps is intentionally conservative: a pattern match contributes
to $\mathcal{C}_{\text{act}}$ only if it appears in executable
or instruction contexts, not in comments or example strings.
This preserves high precision at the cost of recall,
making all IC rates \emph{lower bounds}.

\subsection{Corpus and Results}
\label{sec:meas_results}

We analyze all \NskillsCH skills in \clawhub (March 2026 snapshot).

\begin{table}[t]
\centering
\caption{IC type distribution across \NskillsCH \clawhub skills.
  Pct.\ over skills with $\mathcal{C}_{\text{act}} \neq \varnothing$.}
\label{tab:ic_dist}
\small
\begin{tabular}{llrr}
\toprule
\textbf{Type} & \textbf{Description} & \textbf{Count} & \textbf{\%}\\
\midrule
IC-1 & Full omission  & 13,975 & 75.8\%\\
IC-2 & Partial omission & 2,507 & 13.6\%\\
IC-4 & Bidirectional  &   535 &  2.9\%\\
IC-3 & False claim    &    74 &  0.4\%\\
IC-0 & Accurate       &    74 &  0.4\%\\
\midrule
None & No dangerous cap.\ detected & 1,278 & 6.9\%\\
\bottomrule
\end{tabular}
\end{table}

\textbf{IC-1 dominates}: 75.8\% of skills with dangerous detected
capabilities entirely omit them from descriptions.
Only 0.4\% achieve accurate disclosure.
The mean Capability Recall $\overline{\text{CR}} = \MeanCR$ indicates
that descriptions mention fewer than one in twenty-three
detected capability signals.
Mean CDI $= \MeanCDI$ confirms that nearly all skills have at least
one misrepresented capability family.

\paragraph{Capability family breakdown.}
The most commonly hidden capability families are:
\texttt{SHELL\_EXEC} (hidden in 38.2\% of skills with the capability),
\texttt{FILE\_WRITE} (29.1\%), and \texttt{NETWORK\_EXFIL} (21.4\%).
These three families account for 72\% of the CAPE-validated dangerous sinks (\S\ref{sec:cape}).

\paragraph{Confirmed dangerous sinks.}
From the IC-1 and IC-2 populations, we identify 329 candidate
dangerous sink skills via LSCPS severity scoring.
Human verification of the top 125 yields the \NpairsTotal confirmed
sinks used in \cape (\S\ref{sec:cape}).
Verification assessed whether the detected capability was
(a) actually present in the implementation and
(b) triggerable by agent-side tool invocation without out-of-band
machine access.
To assess labeling reliability, two annotators independently verified
a stratified random sample of 40 sinks (32\% of the verified set).
Inter-annotator agreement (Cohen's $\kappa$) was \textbf{0.87} on
binary dangerous/not-dangerous classification,
indicating near-perfect agreement and supporting the reliability
of the full 125-sink ground truth set.

\subsection{Cross-Platform Validation}

To confirm findings generalize beyond \clawhub,
we analyze 44 skills from skills.sh---an independent, Europe-based
LLM skill marketplace with a different operator and user base---using
the same \lscps pipeline with no parameter changes.
40/44 (90.9\%) exhibit IC-1 or IC-2,
consistent with our primary corpus.
The cross-platform consistency rules out ClawHub-specific effects
and suggests the transparency gap is a structural property of
current skill marketplace designs, not an artifact of any single platform's policies.
Note that this paper's quantitative claims (IC-1 = 75.8\%, CR = 0.042)
are scoped to \clawhub; cross-platform rates serve as directional confirmation only.
